\chapter{Sleep이 먹는 것: data curation, augmentation, replay}
\label{bg:chapter-data-replay}

Sleep-time compute가 training이라면 첫 질문은 ``어떤 optimizer를 쓸까''가 아니라
``어떤 경험을 학습 가능한 증거로 만들까''다. Wake log에는 사용자 발화, model response,
tool result, correction, 반복 조회, 실패 trace가 섞여 있다. 이들을 그대로 next-token data로
넣으면 개인 식별 정보, 틀린 model output, 중복된 쉬운 예, 시간 순서가 깨진 사실이 모두 같은
가중치로 학습된다. 따라서 sleep pipeline의 품질 상한은 대개 data plane에서 정해진다.

\section{Log와 training example은 다르다}

\concept{data-curation}{데이터 큐레이션}은 후보 경험을 수집하고 정규화하며 중복을 제거하고,
목표·권리·민감도·시간 범주를 부여하는 전 과정이다. \concept{data-filtering}{데이터 필터링}은
그중 정해진 기준으로 남길 것과 제외할 것을 가르는 단계다. 둘을 분리해야 ``수집은 했지만
학습에는 쓰지 않음'', ``외부 기억에는 보존하지만 parametric update에는 금지'' 같은 결정을
표현할 수 있다.

한 sleep example $e_i$를 최소한 다음 record로 본다.

\begin{equation}
e_i=(x_i, y_i, t_i, u_i, s_i, q_i, r_i, p_i, v_i),
\end{equation}

여기서 $x_i$는 입력 context, $y_i$는 목표 또는 preference, $t_i$는 event time, $u_i$는
user/tenant scope, $s_i$는 source type, $q_i$는 quality와 uncertainty, $r_i$는 retention/rights,
$p_i$는 provenance pointer, $v_i$는 생성한 model·tool version이다. 실제 시스템에서는 원문과
이 record를 분리해 저장할 수 있다. 원문 삭제 요청을 처리하려면 record가 어느 artifact와
update에 기여했는지도 이어져야 한다.

\concept{provenance}{출처 추적성}은 단순 URL 열이 아니다. ``원본 event $\rightarrow$
정제 record $\rightarrow$ replay batch $\rightarrow$ checkpoint $\rightarrow$ serving version''을
역추적할 수 있어야 한다. provenance가 없으면 잘못된 기억을 발견해도 external index 한 row만
지우면 되는지, LoRA를 폐기해야 하는지, base model을 다시 만들어야 하는지 판단할 수 없다.

\begin{table}[htbp]
\centering
\caption{Wake event를 sleep data로 바꿀 때 필요한 gate}
\label{tab:bg-data-gates}
\small
\begin{tabularx}{\textwidth}{p{0.16\textwidth} X X}
\toprule
Gate & 묻는 질문 & 실패 시 기본 동작 \\
\midrule
scope & 누구의 기억이며 어느 tenant에서만 유효한가? & global 학습 금지, 격리 저장 \\
grounding & correction이나 fact가 신뢰 가능한 source에 연결되는가? & 외부 후보로 격하 또는 제외 \\
temporal & 언제부터 언제까지 참이며 newer record와 충돌하는가? & time interval과 supersession edge 기록 \\
rights/privacy & 저장·재생·가중치 반영·장기 보존 권한이 각각 있는가? & 허용된 destination만 사용 \\
difficulty & 모델이 이미 맞히는 중복 easy example인가? & down-sample 또는 대표 예만 보존 \\
contamination & validation/test 또는 비공개 benchmark 정답이 섞였는가? & quarantine, 평가 집합 재구성 \\
\bottomrule
\end{tabularx}
\end{table}

\section{Augmentation은 의미를 보존한다는 가정이다}

\concept{data-augmentation}{데이터 증강}은 한 예시를 여러 변형으로 늘리는 기법이지만,
핵심은 ``이 변환 아래 목표가 변하지 않는다''는 invariance 가정이다. 이미지 회전처럼 명확한
경우도 있지만 memory data에서는 더 위험하다. 발화를 paraphrase하면 표현 다양성은 늘지만
negation, 날짜, actor, uncertainty가 바뀔 수 있다. context를 잘라내면 정답은 남아 보여도
근거와 privacy scope가 사라질 수 있다.

Sleep용 augmentation은 세 층으로 나누면 안전하다.

\begin{enumerate}
  \item \textbf{surface augmentation}: 표현·순서·format을 바꾸되 entity, time, polarity,
        provenance를 보존한다.
  \item \textbf{contrast construction}: 헷갈리는 old/new fact, 같은 이름의 다른 사람,
        허용/금지 scope를 쌍으로 만들어 decision boundary를 학습한다.
  \item \textbf{counterfactual probe}: 목표 자체를 학습시키기보다 모델이 evidence 없이
        외운 답을 내는지 평가한다. training set과 validation set을 섞지 않는다.
\end{enumerate}

\concept{synthetic-data}{합성 데이터}는 teacher model, simulator, rule로 만든 예시다. 긴
episode를 question--answer pair로 변환하거나 드문 failure를 증폭하는 데 유용하지만, 생성자의
오류와 표현 편향도 함께 복제한다. recursive distillation에서 생성 data가 다시 다음 cycle의
teacher data가 되는 깊이를 record하고, real-evidence anchor를 일정 비율 유지해야 하는 이유가
여기에 있다 \parencite{shumailov2024collapse}.

\section{Replay: 과거를 다시 보여 주는 가장 직접적인 기준선}

\concept{replay-buffer}{리플레이 버퍼}는 과거 경험 일부를 다시 학습시키기 위한 제한된 저장소다.
새 batch $B_{new}$만 학습하지 않고 replay sample $B_{old}$를 섞으면

\begin{equation}
\mathcal{L}_{sleep}=\lambda_{new}\mathcal{L}(B_{new})+
\lambda_{old}\mathcal{L}(B_{old})+lambda_{safe}\mathcal{L}(B_{safety})
\end{equation}

처럼 plasticity와 retention을 직접 trade-off할 수 있다. 작은 episodic memory도 forgetting을
줄이는 강한 baseline이 될 수 있으며 \parencite{chaudhry2019tiny}, class prototype에 가까운
대표 예를 유지하는 방식도 쓰인다 \parencite{rebuffi2017icarl}. 그러나 buffer가 raw episode를
보존하면 저장·privacy 비용이 크고, random sampling은 다수 사용자의 반복 패턴이 희귀하지만
중요한 경험을 밀어낸다.

\concept{coreset}{코어셋}은 전체 데이터의 loss나 representation을 작은 weighted subset으로
근사하려는 관점이다. Sleep system에서는 단순 diversity 외에도 다음 utility를 함께 본다.

\begin{equation}
U(e)=\alpha N(e)+\beta E(e)+\gamma F(e)+\delta R(e)
-\eta C_{store}(e)-\zeta C_{risk}(e),
\end{equation}

$N$은 novelty, $E$는 evidence quality, $F$는 예상 future reuse, $R$은 forgetting 방지 기여다.
이 점수는 established formula가 아니라 controller를 설계하기 위한 decomposition이다. 각 항의
scale과 label delay가 다르므로 하나의 숫자로 만들기 전에 raw features와 의사결정 이유를 남긴다.

\concept{generated-replay}{생성 리플레이}는 raw 과거를 전부 저장하는 대신 generator가 과거
분포를 복원하도록 한다 \parencite{shin2017dgr}. 저장 bytes를 줄일 수 있지만 generator 자체도
잊고, 원본 provenance를 잃으며, tail을 평탄화할 수 있다. 따라서 ``raw replay 대 generated
replay''는 어느 하나를 고르는 문제가 아니라 민감한 원문, 검증된 canonical fact, 일반 skill
pattern에 서로 다른 retention을 주는 tiering 문제다.

\begin{keypoint}[Sleep dataset의 최소 산출물]
Sleep job 입력은 folder 하나가 아니라 immutable dataset manifest여야 한다. manifest에는 source
snapshot, filter version, included/excluded count와 reason, tenant/time scope, dedup cluster,
augmentation lineage, train/validation split, hash, expiry가 포함된다. 같은 checkpoint를 다시 만들 수
없다면 rollback도 감사도 불가능하다.
\end{keypoint}

\section{평가 누수와 시간 누수}

개인화 memory에서는 random split이 특히 위험하다. 같은 episode의 paraphrase가 train과 validation에
나뉘면 memorization을 generalization으로 오해한다. 미래 correction이 과거 시점 평가 입력에 들어가면
temporal leakage가 생긴다. 권장 split unit은 token이 아니라 user, episode, entity thread, time window다.
평가는 ``현재 사실 회상'', ``과거 시점 질의'', ``다른 user로의 비누출'', ``삭제 후 비회상''을 분리한다.

Sleep-time compute에서 data volume은 공짜 scaling 축이 아니다. 중복 experience를 더 많이 replay하면
FLOPs는 늘지만 effective evidence는 거의 늘지 않는다. 그래서 future scaling law의 독립변수는 raw
token 수보다 provenance-weighted novel evidence, replay diversity, recursive depth, validation coverage에
가까워야 한다.
